\section{Related Work}
\label{sec:related-work}

\subsection{Interactive Video Retrieval at VBS}

Interactive systems at VBS combine retrieval models with mechanisms for
iteratively reformulating queries, inspecting candidates, and refining result
sets. VISIONE supports textual, visual-similarity, object, color, and temporal
search, with subsequent versions updating its retrieval models and result
presentation~\cite{amato2022visione,amato2023visione,amato2024visione}.
Vibro emphasizes efficient browsing together with rich-text and temporal
search~\cite{hezel2022vibro}, while PraK combines text-query reformulation,
example-image search, filtering, and relevance
feedback~\cite{lokoc2024prak}. Other systems emphasize complementary forms of
interaction: diveXplore provides browsing and exploration views over clustered
results~\cite{schoeffmann2024divexplore}, vitrivr offers a modular retrieval
engine for composing query-processing mechanisms~\cite{gasser2024vitrivr},
and Exquisitor combines conversational query refinement with user relevance
feedback~\cite{khan2024exquisitor,sharma2025exquisitor}.

Evidence from VBS evaluations further shows that effective video search depends
not only on the initial retrieval model, but also on subsequent query
reformulation, browsing, and result inspection~\cite{schall2024interactive,jackl2026interaction}. In a post-evaluation of VBS 2022, VISIONE achieved the
best text-to-image retrieval performance among the three evaluated systems,
whereas Vibro obtained the highest overall result; browsing and image-based
retrieval were identified as important factors in its
performance~\cite{schall2024interactive}. More recently, an interaction-log
analysis of two systems at VBS 2026 found iterative text-query reformulation
followed by result-set inspection to be the dominant strategy across most task
categories~\cite{jackl2026interaction}. Together, these observations motivate
interactive retrieval workflows in which intermediate search states can be
revised rather than treating the initial ranked list as the final system
output~\cite{schall2024interactive,jackl2026interaction}. SHI follows this
perspective while exposing two additional interaction states: the event-level
interpretation of the query and the temporal alignment within a retrieved
video.

\subsection{Query Reformulation and Temporal Interaction}

Interactive query reformulation has primarily operated on the query itself or
on retrieved candidates. Loko{\v{c}} et al.\ derive class suggestions from
intermediate CLIP results to assist users in revising textual
queries~\cite{lokoc2023queryreformulation}. Exquisitor instead combines
natural-language refinement with relevance feedback, with its later version
integrating both mechanisms into a unified search
interface~\cite{khan2024exquisitor,sharma2025exquisitor}. vitrivr has also
explored explainability mechanisms intended to support result inspection and
relevance feedback, alongside interaction with temporally contextualized
queries~\cite{heller2021explainable}. In contrast, Query Hypothesis makes the
system's interpretation of a multi-event request an explicit interaction
object: users can inspect the ordered event decomposition and revise its
structure before retrieval.

Temporal structure has likewise been incorporated into interactive video
retrieval. Heller et al.\ support temporal-context queries composed from
multiple retrieval modalities~\cite{heller2022temporalqueries}, while
Exquisitor 2026 introduces sequence-chain retrieval with reciprocal-rank fusion
for multi-event temporal queries~\cite{khan2026temporalqueries}. Beyond VBS,
interactive video-corpus moment retrieval has used user feedback to guide
subsequent navigation and recommendation rounds~\cite{ma2022ivcmr}. These
approaches use interaction to refine queries, rankings, or navigation through
candidate moments. EventTrail instead makes the current multi-event temporal
alignment within a retrieved video directly inspectable and revisable, allowing
local user corrections to update the hypothesis for the complete ordered event
sequence.
